\documentclass[11pt]{article}

\usepackage[margin=1in]{geometry}
\usepackage{amsmath,amssymb,booktabs,longtable,array}
\usepackage[hidelinks]{hyperref}

\newcommand{\E}{\mathbb E}
\newcommand{\R}{\mathbb R}

\title{Notation List and Consistency Guide}
\author{Azevedo--Wolff (working document)}
\date{\today}

\begin{document}
\maketitle

\section{Notation policy}

This document records the notation currently used in the paper, regular
appendix, and online supplementary appendix.  The guiding rules are:
\begin{enumerate}
    \item A symbol has one substantive role throughout the project whenever
    possible.
    \item Lowercase $w$ is payment (the wage), while uppercase $W$ is final
    wealth.  Utility is always written as a function of final wealth, $u(W)$.
    \item $L$ is reserved for the \emph{wage link function}; $\ell$ is
    reserved for the \emph{utility link function}.  A subscript $0$ denotes
    the corresponding nonbinding/base link function.
    \item In the body, ``wage link function'' is used when $L$ is introduced
    or discussed.  The proofs appendix states once that only utility link
    functions are used there and then says simply ``link function'' for
    brevity.
    \item Uppercase $S(y\mid a)$ is the action score at an arbitrary action;
    lowercase $s(y)$ is the score after evaluation at the intended action
    $a_0$.
    \item A vertical bar is reserved for conditioning, as in $f(y\mid a)$.
    Parameters of contracts are therefore subscripts, as in
    $w_{\lambda,\mu}$.
\end{enumerate}

The notation $u(W)$ in the closed-form table is intentional and consistent:
$W$ is total wealth.  The table now makes explicit that the wage link function $L$
returns canonical total wealth and that the payment is obtained by subtracting
$w_0$.

\section{Economic primitives and contract problem}

\begin{longtable}{>{$}l<{$} p{0.77\textwidth}}
\toprule
\text{Symbol} & \text{Definition and use} \\
\midrule
\endhead
\mathcal A & Compact action set, $\mathcal A\subseteq\R_+$. \\
a & Agent's action. \\
a_0 & Fixed interior intended action. \\
y & Realized output. \\
Y & Output random variable, used when expectation notation is convenient. \\
\mathcal Y & Common output support/state space. \\
\nu & Common dominating measure.  Integrals cover densities and probability
masses; $\nu$ is counting measure in discrete examples. \\
f(y\mid a) & Density or probability mass function of output conditional on
action.  It is strictly positive on the common support. \\
f_a,\ f_{aa} & First and second derivatives of $f(y\mid a)$ with respect to
action. \\
c(a) & Agent's action cost. \\
\underline c'' & Uniform lower curvature of cost,
$\inf_{a\in\mathcal A}c''(a)$. \\
w_0 & Agent's initial wealth. \\
w(y) & Nonnegative contractual payment (wage) after output $y$. \\
\mathcal W & Set of feasible payment contracts. \\
W & Final wealth, with $W=w_0+w(y)$ when induced by a contract. \\
u(W) & Bernoulli utility over final wealth. \\
\bar U & Reservation utility.  Uppercase $U$ distinguishes this economic
utility level from the utility ceiling $\bar u$. \\
\bar w & Certainty-equivalent reservation wage, used in the motivating
figures to express $\bar U$ in monetary units. \\
U(w,a) & Agent's expected utility net of action cost under payment contract
$w$: $\int u(w_0+w(y))f(y\mid a)d\nu(y)-c(a)$. \\
C(w,a) & Principal's expected compensation cost:
$\int w(y)f(y\mid a)d\nu(y)$. \\
w^*_{\bar U} & Relaxed-optimal (and, when the FOA is valid, globally optimal)
payment contract at reservation utility $\bar U$. \\
\mathcal L & Lagrangian in payment-contract coordinates.  Calligraphic
$\mathcal L$ is not the wage link function $L$. \\
\text{IR} & Individual rationality, $U(w,a_0)\geq\bar U$. \\
\text{GIC} & Global incentive compatibility,
$U(w,a_0)\geq U(w,a)$ for all $a\in\mathcal A$. \\
\text{LIC} & Local incentive compatibility, $U_a(w,a_0)=0$. \\
\bottomrule
\end{longtable}

\section{Scores and distribution notation}

\begin{longtable}{>{$}l<{$} p{0.77\textwidth}}
\toprule
\text{Symbol} & \text{Definition and use} \\
\midrule
\endhead
S(y\mid a) & Action score at arbitrary $a$:
$\partial_a\log f(y\mid a)=f_a(y\mid a)/f(y\mid a)$. \\
s(y) & Score at the intended action, $S(y\mid a_0)$.  It has mean zero
under $f(\cdot\mid a_0)$. \\
\sigma & Scale or standard deviation in parametric distributions. \\
\sigma_s^2 & Second moment of the intended-action score,
$\int s(y)^2f(y\mid a_0)d\nu(y)$. \\
h(y) & Base density/mass in an exponential family, or the density of a shock
in additive and multiplicative models. \\
\eta(a) & Natural parameter in an exponential family. \\
T(y) & Sufficient statistic in an exponential family. \\
\psi(a) & Log-normalizing term in an exponential family.  The symbol $\psi$
is used rather than $A$ to avoid collision with absolute risk aversion
$A(W)$. \\
n & Number of binomial trials or fixed gamma shape; the context identifies
the distribution. \\
d & Degrees of freedom for Student's $t$.  This avoids overloading the
dominating measure $\nu$. \\
X & Exogenous shock in additive or multiplicative output models. \\
\E_a & Expectation under $f(\cdot\mid a)$. \\
y_-,y_0,y_+ & Low, middle, and high states in the supplementary discrete
counterexample.  These replace $L,M,H$, reserving $L$ for the wage link function. \\
\bottomrule
\end{longtable}

\section{Link functions and canonical contracts}

Let $z$ denote the scalar index to which a link function is applied.

\begin{longtable}{>{$}l<{$} p{0.77\textwidth}}
\toprule
\text{Symbol} & \text{Definition and use} \\
\midrule
\endhead
m_0 & Marginal-cost threshold at zero payment,
$m_0=1/u'(w_0)$. \\
L_0(z) & Base (nonbinding) wage link function,
$L_0(z)=(u')^{-1}(1/z)$.  It returns total wealth. \\
L(z) & Limited-liability wage link function,
$L(z)=L_0(z\vee m_0)$.  The canonical payment is therefore $L(z)-w_0$. \\
\ell_0(z) & Base utility link function,
$\ell_0(z)=u(L_0(z))=(k')^{-1}(z)$. \\
\ell(z) & Limited-liability utility link function,
$\ell(z)=\ell_0(z\vee m_0)=u(L(z))$. \\
k & Inverse utility, $k=u^{-1}$. \\
\lambda & Lagrange multiplier on individual rationality; also the parameter
indexing the locally incentive-compatible Pareto path. \\
\mu & Lagrange multiplier on local incentive compatibility. \\
w_{\lambda,\mu}(y) & Canonical payment contract,
$L(\lambda+\mu s(y))-w_0$. \\
W_{\lambda,\mu}(y) & Canonical final wealth,
$L(\lambda+\mu s(y))$; used in the supplementary link-function derivations. \\
v_{\lambda,\mu}(y) & Canonical utility contract,
$\ell(\lambda+\mu s(y))$. \\
\widetilde\mu(\lambda) & Unique positive LIC multiplier for a fixed
$\lambda$. \\
v_\lambda & LIC-path contract
$v_{\lambda,\widetilde\mu(\lambda)}$. \\
I(\lambda,\mu) & Local-incentive functional,
$\int\ell(\lambda+\mu s(y))s(y)f(y\mid a_0)d\nu(y)$. \\
\rho & Finite asymptotic link-function coefficient
$\lim_{z\to\infty}z\ell_0'(z)$ in the safe-tail bound.  This was changed from
$L$ so that $L$ remains reserved for the wage link function. \\
\bottomrule
\end{longtable}

For the utility specifications used in the paper:
\[
\begin{array}{c|c|c}
 u(W) & L(z) & \ell(z)\\ \hline
 \log W & z\vee w_0 & \log(z\vee w_0)\\[2pt]
 \dfrac{W^{1-\gamma}}{1-\gamma}
 & (z\vee w_0^\gamma)^{1/\gamma}
 & \dfrac{(z\vee w_0^\gamma)^{(1-\gamma)/\gamma}}{1-\gamma}\\[7pt]
 -\dfrac{e^{-\alpha W}}{\alpha}
 & \dfrac{\log(z\vee e^{\alpha w_0})}{\alpha}
 & -\dfrac{1}{\alpha(z\vee e^{\alpha w_0})}
\end{array}
\]
Here $\gamma$ is relative risk aversion for CRRA and $\alpha$ is absolute risk
aversion for CARA.

\section{Utility-contract coordinates used in the proofs}

\begin{longtable}{>{$}l<{$} p{0.77\textwidth}}
\toprule
\text{Symbol} & \text{Definition and use} \\
\midrule
\endhead
\underline u & Limited-liability utility floor, $u(w_0)$. \\
\bar u & Utility ceiling, $\lim_{W\to\infty}u(W)$, possibly infinite. \\
\Delta u & Available utility range, $\bar u-\underline u$. \\
v(y) & Utility contract, $u(w_0+w(y))$. \\
\mathcal V & Feasible utility-contract set. \\
\widehat U(v,a) & Expected utility in utility-contract coordinates,
$\int v(y)f(y\mid a)d\nu(y)-c(a)$. \\
\widehat C(v,a) & Expected payment cost in utility-contract coordinates,
$\int[k(v(y))-w_0]f(y\mid a)d\nu(y)$. \\
\widehat{\mathcal L} & Lagrangian in utility-contract coordinates. \\
\widetilde U(\lambda) & Utility attained along the LIC path,
$\widehat U(v_\lambda,a_0)$. \\
\widetilde C(\lambda) & Compensation cost along the LIC path,
$\widehat C(v_\lambda,a_0)$. \\
U_L & Lower reservation-utility frontier, $\widetilde U(0)$.  The subscript
means ``lower'' and is not a use of the standalone wage-link-function symbol $L$. \\
U_R & Supremum feasible reservation utility for the relaxed problem. \\
U^* & Reservation-utility threshold above which the main theorem validates
the first-order approach. \\
C^*(\bar U) & Optimal expected compensation cost at reservation utility
$\bar U$. \\
\Lambda(\bar U) & Set of IR multipliers associated with reservation utility
$\bar U$. \\
\bottomrule
\end{longtable}

\section{Safe low-score notation}

\begin{longtable}{>{$}l<{$} p{0.77\textwidth}}
\toprule
\text{Symbol} & \text{Definition and use} \\
\midrule
\endhead
q & Generic nonpositive score cutoff. \\
\mathcal Q_{\mathrm{safe}} & Cutoffs $q\leq0$ for which
$f_{aa}(y\mid a)\geq0$ for every action on $\{s(y)\leq q\}$. \\
\bar q_{\mathrm{safe}} & Supremum of safe cutoffs. \\
\mathrm{Cap}(q) & Cutoff incentive capacity,
$-\int_{\{s(y)\leq q\}}s(y)f(y\mid a_0)d\nu(y)$. \\
\mathrm{Curv}(q) & Positive curvature above cutoff $q$:
$\sup_{a\in\mathcal A}\int_{\{s(y)>q\}}[s(y)f_{aa}(y\mid a)]_+d\nu(y)$. \\
\mathrm{Cap}_{\mathrm{safe}} & Aggregate safe incentive capacity, expressed
in the body as the limiting safe-tail integral. \\
\mathrm{Curv}_{\mathrm{safe}} & Infimum of $\mathrm{Curv}(q)$ over safe
cutoffs. \\
\underline s & Essential infimum of $s(y)$ under the intended-action output
distribution. \\
\mathrm{Width}_{\mathrm{safe}} & Maximum score width covered by a safe
low-score region. \\
q_0,q_1 & Two safe cutoffs in the proof, with $q_0<q_1\leq0$; $q_0$ controls
kink placement and $q_1$ controls positive curvature. \\
q^{\mathrm{kink}}(\lambda) & Limited-liability score cutoff along the LIC
path, $(m_0-\lambda)/\widetilde\mu(\lambda)$. \\
\mu_q(\lambda) & Trial multiplier placing the kink at $q$,
$(m_0-\lambda)/q$. \\
\delta_\lambda(t) & Centered utility-link-function increment,
$\ell(\lambda+\widetilde\mu(\lambda)t)-\ell(\lambda)$. \\
\lbrack x\rbrack_+ & Positive part, $\max\{x,0\}$. \\
x\vee y & Maximum of $x$ and $y$. \\
\bottomrule
\end{longtable}

\section{Risk aversion and supplementary asymptotics}

\begin{longtable}{>{$}l<{$} p{0.77\textwidth}}
\toprule
\text{Symbol} & \text{Definition and use} \\
\midrule
\endhead
A(W) & Absolute risk aversion, $-u''(W)/u'(W)$. \\
P(W) & Absolute prudence, $-u'''(W)/u''(W)$. \\
\alpha & CARA coefficient. \\
\gamma & CRRA coefficient.  Log utility is the limiting $\gamma=1$ case. \\
p & Transformed CRRA exponent in the low-risk-aversion supplement,
$p=(1-\gamma)/\gamma$. \\
r & Moment order in the asymptotic-affinity argument.  This avoids collision
with the limited-liability threshold $m_0$. \\
\theta & Positive scaling argument used to bracket the LIC multiplier. \\
T_\lambda & Score truncation threshold in the asymptotic-affinity proof. \\
K & Generic positive finite bound, allowed to change between displays. \\
\widehat a_i & Candidate global-deviation action in the numerical active set. \\
\widehat\mu_i & Multiplier on the global incentive constraint at
$\widehat a_i$. \\
\boldsymbol{\widehat a},\boldsymbol{\widehat\mu} & Vectors of active
deviation actions and their multipliers. \\
\mathcal D & Lagrange dual function in the numerical method. \\
\mathbf y_{\mathrm{grid}} & Numerical output grid. \\
\bottomrule
\end{longtable}

\section{Audit decisions and remaining cautions}

The current review made the following consistency changes in the manuscript:
\begin{itemize}
    \item standardized the names \emph{wage link function} $L$ and
    \emph{utility link function} $\ell$;
    \item stated the proofs-appendix convention that ``link function'' alone
    means utility link function;
    \item changed the discrete counterexample states from $L,M,H$ to
    $y_-,y_0,y_+$;
    \item changed the safe-tail asymptotic constant from $L$ to $\rho$;
    \item changed Student's-$t$ degrees of freedom from $\nu$ to $d$, leaving
    $\nu$ exclusively as the dominating measure;
    \item changed the exponential-family log normalizer from $A(a)$ to
    $\psi(a)$, leaving $A(W)$ for absolute risk aversion;
    \item changed the supplementary moment order from $m$ to $r$, avoiding
    proximity to the threshold $m_0$; and
    \item removed the introductory example's old payment-domain notation
    $u(x)=\log(w_0+x)$ in favor of utility over final wealth, $u(W)=\log W$.
\end{itemize}

Two benign context-dependent reuses remain.  The symbol $n$ denotes binomial
trials or gamma shape in separate distribution rows.  The letter $L$ appears as
a subscript in $U_L$ meaning ``lower frontier,'' but never as a standalone
state or scalar.  Calligraphic $\mathcal L$ remains the conventional
Lagrangian and is visually distinct from the wage link function $L$.

When adding material, do not use standalone $L$ for a state, limit, loss, or
likelihood; do not use $\ell$ for anything other than the utility link function; and do
not revert to $u(w)$ or $u(x)$ when utility is over final wealth.

\end{document}
